\chapter{Communication Protocol Specification}

\section{Protocol Overview}

This chapter defines the structured packet-based communication protocol used for bidirectional data exchange between the embedded firmware and the TeleCon4ESP32 Android application.

The protocol is designed to achieve:

\begin{itemize}
	\item Deterministic packet parsing
	\item Robust synchronization recovery
	\item Low processing overhead
	\item Fixed-length command handling
	\item Structured telemetry formatting
\end{itemize}

All communication occurs over a Bluetooth serial transport layer. The protocol operates independently of the underlying physical transport once a byte stream is established.

% ------------------------------------------------

\section{Communication Model}

The communication model is packet-oriented.  
All data exchanged between the application and firmware is encapsulated within structured packets.

Two packet categories exist:

\begin{itemize}
	\item Control Packets (Downlink)
	\item Telemetry Packets (Uplink)
\end{itemize}

The firmware parser continuously monitors the incoming byte stream and attempts to detect valid packet frames.

% ------------------------------------------------

\section{Packet Framing Strategy}

To guarantee reliable synchronization within a continuous byte stream, each packet contains:

\begin{itemize}
	\item A Start-of-Frame (SOF) marker
	\item A Payload section
	\item A Checksum field
\end{itemize}

\subsection{General Packet Structure}

\begin{center}
	\begin{tabular}{|c|c|c|}
		\hline
		SOF & Payload & Checksum \\
		\hline
	\end{tabular}
\end{center}

\subsection{Start-of-Frame Marker}

The SOF marker is a predefined byte sequence used to identify the beginning of a valid packet.  

Its purpose is to:

\begin{itemize}
	\item Allow parser resynchronization after byte loss
	\item Prevent payload misinterpretation
	\item Reduce false-positive packet detection
\end{itemize}

The SOF pattern is unique and not permitted to appear as a valid packet prefix.

% ------------------------------------------------

\section{Control Packet Format}

Control packets are transmitted from the TeleCon4ESP32 application to the firmware.

\subsection{Control Packet Structure}

\begin{center}
	\begin{tabular}{|c|c|c|c|}
		\hline
		SOF & Control Fields & Reserved & Checksum \\
		\hline
	\end{tabular}
\end{center}

\subsection{Control Fields}

The control payload contains fixed-size fields representing actuator commands and output states.

Typical control fields may include:

\begin{itemize}
	\item Motor speed values
	\item Steering position
	\item Auxiliary servo position
	\item Output enable flags
\end{itemize}

Each field has:

\begin{itemize}
	\item Defined byte width
	\item Defined numerical range
	\item Defined scaling factor
\end{itemize}

All control packets are fixed-length to simplify parsing logic and maintain constant processing time.

% ------------------------------------------------

\section{Telemetry Packet Format}

Telemetry packets are transmitted from the firmware to the TeleCon4ESP32 application.

\subsection{Telemetry Packet Structure}

\begin{center}
	\begin{tabular}{|c|c|c|c|}
		\hline
		SOF & Telemetry Fields & Status Flags & Checksum \\
		\hline
	\end{tabular}
\end{center}

\subsection{Telemetry Fields}

Telemetry payload may include:

\begin{itemize}
	\item Sensor measurements
	\item Actuator state feedback
	\item Input states
	\item System indicators
\end{itemize}

Telemetry structure is deterministic and ordered to ensure consistent parsing on the application side.

% ------------------------------------------------

\section{Checksum and Error Detection}

Each packet terminates with a checksum byte.

\subsection{Checksum Purpose}

The checksum enables:

\begin{itemize}
	\item Detection of corrupted packets
	\item Rejection of malformed data
	\item Prevention of unsafe actuator updates
\end{itemize}

\subsection{Checksum Model}

The checksum is computed as a deterministic function of all payload bytes:

\[
Checksum = f(b_1, b_2, ..., b_n)
\]

Where:

\begin{itemize}
	\item $b_i$ represents payload bytes
	\item $f$ is a defined aggregation function (e.g., modulo sum)
\end{itemize}

Packets failing checksum validation are discarded without affecting actuator states.

% ------------------------------------------------

\section{Parser Synchronization Model}

The parser operates as a state machine with the following states:

\begin{enumerate}
	\item Idle (waiting for SOF)
	\item Receiving Payload
	\item Validating Checksum
	\item Packet Accepted / Rejected
\end{enumerate}

If checksum validation fails, the parser returns to the Idle state and resumes searching for the next SOF marker.

This ensures recovery from:

\begin{itemize}
	\item Byte loss
	\item Transmission noise
	\item Partial packet corruption
\end{itemize}

% ------------------------------------------------

\section{Timing Characteristics}

Control packets are processed within a single control loop iteration once fully received and validated.

Telemetry packets are generated at a deterministic rate defined by the control loop frequency.

The protocol avoids variable-length parsing to ensure bounded processing time.

% ------------------------------------------------

\section{Protocol Determinism}

The protocol enforces determinism through:

\begin{itemize}
	\item Fixed packet length
	\item Defined byte ordering
	\item Consistent field offsets
	\item Stateless packet interpretation
\end{itemize}

This guarantees constant-time decoding complexity.

% ------------------------------------------------

\section{Extensibility Strategy}

Protocol extensibility is achieved through:

\begin{itemize}
	\item Reserved payload bytes
	\item Version identifier field (optional)
	\item Backward-compatible field extension
\end{itemize}

Future fields may be appended without altering parser stability if packet size constraints are maintained.

% ------------------------------------------------

\section{Protocol Robustness Considerations}

The protocol design accounts for:

\begin{itemize}
	\item Bluetooth latency variability
	\item Occasional packet loss
	\item Transmission jitter
\end{itemize}

No partial command execution occurs.  
Only fully validated packets update control variables.

This prevents unintended actuator behavior.

% ------------------------------------------------

\section{Summary}

The communication protocol provides:

\begin{itemize}
	\item Structured bidirectional packet exchange
	\item Deterministic decoding behavior
	\item Error detection via checksum validation
	\item Synchronization recovery capability
	\item Fixed processing complexity
\end{itemize}

The following chapter details how this protocol is implemented within the firmware module structure.